\documentclass[10pt,reqno]{amsart}

\usepackage[T1]{fontenc}
\usepackage[utf8]{inputenc}
\usepackage{newtxtext,newtxmath}
\usepackage{microtype}
\usepackage{mathtools}
\usepackage{enumitem}
\usepackage[nameinlink,noabbrev]{cleveref}
\usepackage[a4paper,margin=28mm,headheight=15pt]{geometry}

\setlength{\parskip}{2pt}
\setlist[enumerate]{leftmargin=7mm,itemsep=2pt,topsep=3pt}
\emergencystretch=2em

\newtheorem{theorem}{Theorem}[section]
\newtheorem{proposition}[theorem]{Proposition}
\newtheorem{lemma}[theorem]{Lemma}
\newtheorem{corollary}[theorem]{Corollary}
\theoremstyle{definition}
\newtheorem{definition}[theorem]{Definition}
\newtheorem{remark}[theorem]{Remark}

\title[A topological obstruction to an inverse arboreal problem]
{A Topological Obstruction to the Literal Inverse Arboreal Galois Problem}

\author{Argus Open Research Project}
\address{Argus Open Research Project}
\date{August 30, 2026}

\subjclass[2020]{Primary 12F10; Secondary 20E08, 37P15}
\keywords{arboreal Galois representation, rooted tree, profinite group, nonclosed subgroup}

\begin{document}

\begin{abstract}
An inverse problem asks whether an arbitrary subgroup of the automorphism
group of a regular rooted tree can occur as the image of an arboreal Galois
representation.  Taken literally, the answer is negative for a purely
topological reason.  The image of a continuous homomorphism from an absolute
Galois group into the profinite automorphism group of a rooted tree is closed.
On the other hand, the finitary automorphism group of the infinite rooted
binary tree is a dense proper subgroup and hence is not closed.  This gives a
counterexample to the unrestricted arbitrary-subgroup formulation.  It does
not address the inverse problem for prescribed closed subgroups.
\end{abstract}

\maketitle

\section{The inverse problem and its topology}

Let \(T_d\) be the infinite regular rooted \(d\)-ary tree.  An arboreal
Galois representation associated with a polynomial and a basepoint is a
continuous homomorphism
\[
 \rho:\operatorname{Gal}(K^{\mathrm s}/K)
 \longrightarrow \operatorname{Aut}(T_d).
\]
The absolute Galois group carries its Krull topology.  The tree automorphism
group carries the inverse-limit topology
\[
 \operatorname{Aut}(T_d)
 \simeq \varprojlim_n \operatorname{Aut}(T_{d,n}),
\]
where \(T_{d,n}\) is the finite truncation through level \(n\).

The literal form of the inverse problem asks whether every abstract subgroup
of \(\operatorname{Aut}(T_d)\) can arise as such an image.  The following
standard observation already imposes a necessary condition.

\begin{proposition}\label{prop:closed}
The image of every arboreal Galois representation is a closed subgroup of
\(\operatorname{Aut}(T_d)\).
\end{proposition}

\begin{proof}
The Krull-topological absolute Galois group is profinite and therefore
compact.  The group \(\operatorname{Aut}(T_d)\), being an inverse limit of
finite discrete groups, is Hausdorff.  The continuous image of a compact
space is compact, and compact subsets of a Hausdorff space are closed.
\end{proof}

\section{The finitary subgroup}

It remains to exhibit an abstract subgroup that violates the necessary
closedness condition.  We work with the binary tree
\[
 T=\{0,1\}^{*},
\]
whose root is the empty word and whose children below a word \(w\) are
\(w0\) and \(w1\).

For each \(n\geq 0\), let \(T_n\) denote the truncation through level \(n\).
Every automorphism \(h\in\operatorname{Aut}(T_n)\) has a trivial extension
to \(T\): if a word of length greater than \(n\) is written uniquely as
\(uv\), with \(\lvert u\rvert=n\), define
\[
 \iota_n(h)(uv)=h(u)v.
\]
Define
\[
 F=\bigcup_{n\geq0}\iota_n\bigl(\operatorname{Aut}(T_n)\bigr).
\]
Equivalently, \(F\) consists of those automorphisms whose local actions are
trivial below some finite level.  It is called the finitary automorphism
group.

\begin{lemma}\label{lem:dense}
The subgroup \(F\) is dense in \(\operatorname{Aut}(T)\).
\end{lemma}

\begin{proof}
Let
\[
 \pi_m:\operatorname{Aut}(T)\longrightarrow\operatorname{Aut}(T_m)
\]
be restriction.  Given \(a\in\operatorname{Aut}(T_m)\), the trivial
extension \(\iota_m(a)\) lies in \(F\) and restricts to \(a\).  Hence
\[
 \pi_m(F)=\operatorname{Aut}(T_m)
 \qquad\text{for every }m.
\]
Every nonempty basic open subset of the inverse-limit topology is determined
by a finite-level restriction.  Since \(F\) realizes every such restriction,
every nonempty basic open set meets \(F\).
\end{proof}

\begin{lemma}\label{lem:proper}
The subgroup \(F\) is a proper subgroup of \(\operatorname{Aut}(T)\).
\end{lemma}

\begin{proof}
Each group \(\operatorname{Aut}(T_n)\) is finite, so \(F\) is countable.
The full automorphism group is uncountable.  Indeed, the vertices
\[
 v_n=1^n0\qquad(n\geq0)
\]
are pairwise prefix-incomparable, and their descendant subtrees are
disjoint.  For each binary sequence \(\epsilon=(\epsilon_n)_{n\geq0}\),
simultaneously interchange the two child subtrees below \(v_n\) precisely
when \(\epsilon_n=1\).  Disjointness of the supports makes this a tree
automorphism, and distinct sequences give distinct automorphisms.  Thus
\(\{0,1\}^{\mathbb N}\) injects into \(\operatorname{Aut}(T)\).
\end{proof}

\begin{corollary}\label{cor:nonclosed}
The finitary subgroup \(F\) is not closed in \(\operatorname{Aut}(T)\).
\end{corollary}

\begin{proof}
By \Cref{lem:dense}, the closure of \(F\) is the full automorphism group.
By \Cref{lem:proper}, \(F\) is not the full automorphism group.
\end{proof}

\section{The counterexample}

\begin{theorem}\label{thm:main}
There is a subgroup of the automorphism group of a regular rooted tree that
cannot be the image of an arboreal Galois representation.  In particular,
the inverse problem with an unrestricted arbitrary-subgroup quantifier is
false.
\end{theorem}

\begin{proof}
Take the finitary subgroup
\[
 F\leq\operatorname{Aut}(T)
\]
of the infinite rooted binary tree.  It is a subgroup by construction and is
not closed by \Cref{cor:nonclosed}.  If \(F\) were the image of an arboreal
Galois representation, \Cref{prop:closed} would make it closed, a
contradiction.
\end{proof}

\begin{remark}[Scope]\label{rem:scope}
The obstruction uses only the omission of a closedness hypothesis.  It gives
no nonrealizability result for any particular closed subgroup of
\(\operatorname{Aut}(T_d)\).  Thus the natural inverse problem for prescribed
closed subgroups remains separate from the literal statement treated here.
\end{remark}

\section{Discussion}

The argument is independent of arithmetic properties of a polynomial.  Its
content is that the topology is part of the data of an arboreal image: an
abstract subgroup that is not closed cannot occur, regardless of the field,
basepoint, or dynamical construction.  Consequently, a well-posed inverse
problem must at least restrict the prescribed group to a closed subgroup of
the profinite tree automorphism group.

\appendix

\section{Checklist for verification}

\begin{enumerate}
  \item Verify that the arboreal representation is continuous for the Krull
  and inverse-limit topologies.
  \item Check compactness of the absolute Galois group and the Hausdorff
  property of \(\operatorname{Aut}(T_d)\).
  \item Verify that every finite-level automorphism has the displayed trivial
  extension and hence that every finite projection of \(F\) is surjective.
  \item Check that \(F\) is countable while \(\operatorname{Aut}(T)\) is
  uncountable.
  \item Conclude that \(F\) is dense, proper, and nonclosed.
  \item Keep the unrestricted abstract-subgroup statement separate from the
  closed-subgroup inverse problem.
\end{enumerate}

\begin{thebibliography}{9}

\bibitem{AHPW}
J.~Anderson, S.~Hamblen, B.~Poonen, and L.~Walton,
\emph{Local arboreal representations},
Int. Math. Res. Not. IMRN (2018), no.~19, 5974--5994;
arXiv:1610.04969v2.

\bibitem{AIM}
American Institute of Mathematics,
\emph{The Galois theory of orbits in arithmetic dynamics},
problem list from the workshop held May 16--20, 2016, Problem~14.1.

\end{thebibliography}

\end{document}
